EdgeTAM, SAM 2'nin video nesne takibi görevini edge cihazlarda
gerçek zamanlı çalıştırabilmek için memory attention modülünü bir 2D
Spatial Perceiver ile sıkıştıran bir mimaridir. Bu çalışma, EdgeTAM'ın
checkpoint'ini NVIDIA Jetson AGX Orin üzerinde çalışan,
kare başına dört ayrı TensorRT motoruna (image encoder, memory
attention, memory encoder, SAM head) fp16 hassasiyetinde ve CUDA Graph
tekrar oynatmasıyla taşımaktadır.

Zamanın çoğunun image encoder'da geçtiği varsayımı ölçümle
çürütülmüştür: kare başına işlem yükünün (FLOP) \%61,9'u memory
attention'da, yalnızca \%26,8'i image encoder'dadır. Bu nedenle
çalışmanın kapsamı tek bir modülü değil, EdgeTAM'ın Python durum takibiyle
(memory bank \texttt{dict}'i) ayrıştığı dört sınırı kapsayacak şekilde
genişletilmiştir.

Cihaz üzerinde ölçülen sonuçlar: uçtan uca inference hızı stok PyTorch'ta
88,30~ms/kare iken TensorRT + CUDA Graph ile 26,40~ms/kareye
inmiş (\textbf{3,34×} hızlanma, 37,88~FPS), kuyruk gecikmesi
(p99)~39,90~ms ile 40~ms hedefinin altında kalmıştır. Sentetik olarak
üretilmiş 500 karelik bir klip boyunca fp16 motorlarının maskeleri fp32 referansına karşı
\textbf{0,9993} ortalama IoU ile eşleşmiş ve hata klip boyunca
birikmemiştir (−0,0012 IoU/1000 kare eğilim). Ayrıca ön-işleme
adımının (kare boyutlandırma, normalize etme) gözden kaçırılmaması
gerektiği; ölçülmediğinde 1024 çözünürlükte modelin kendisi kadar maliyetli
olabildiği, doğru şekilde yeniden yazıldığında ise 4--5× ucuzlatılabildiği
gösterilmiştir.

Rapor; SAM'den SAM~2'ye ve EdgeTAM'a mimari evrimi, dört motorluk ayrışma
kararının gerekçesini, ölçüm metodolojisini ve cihaz üzerinde alınan
tüm sonuçları kapsar. NVIDIA TensorRT Model Optimizer ile gerçek INT8
motoru ve dört motoru tek motora birleştirme sorusu gelecek çalışma
olarak bırakılmıştır.
